\input{./._relpath-to-latexroot.ltx}
\documentclass[\latexroot/\projectname]{subfiles}
\input{\latexroot/@resources/subfile-setup.ltx}

\begin{document}

\section{Introduction}
\whenintegrated{\label{sec:intro}}
\setcounter{page}{0}\pagenumbering{arabic}

Although many empirical studies argue that potential and existing entrepreneurs face borrowing constraints, so far there has been little work on how these constraints affect aggregate capital accumulation and wealth inequality through entrepreneurial choices. Do these financial constraints hamper aggregate capital accumulation and, if so, how big is this effect? What effect do these constraints have on wealth inequality: do they exacerbate it or mitigate it? These are potentially important forces to understand the consequences of policy reforms that affect the tightness of these borrowing constraints, such as changes in the leniency of bankruptcy laws and in the degree of enforcement of property rights.

In this paper we analyze the role of borrowing constraints as determinants of entrepreneurial decisions (entry, continuation, investment, and saving), and their effects on wealth inequality and aggregate capital accumulation, in a framework that matches the observed wealth inequality very closely. In the presence of borrowing constraints, the decision to invest, the fraction of entrepreneurs, and the size distribution of firms depend on the distribution of assets in the economy. Because of this interaction, it is key to perform such an analysis in a model that matches well the extreme concentration of wealth observed in the data.

We find that more restrictive borrowing constraints generate less inequality in wealth holdings but also reduce average firm size, the number of people engaging in entrepreneurial activities, and aggregate capital accumulation. Our results also indicate that voluntary bequests are an important channel allowing some high-ability workers to establish or enlarge an entrepreneurial activity. If there were only accidental bequests, there would be fewer very large firms and less aggregate capital, but also less wealth concentration.

These findings are based on a quantitative life cycle model with altruism across generations and entrepreneurial choice, in an environment in which debt repayment cannot be perfectly enforced. The amount that entrepreneurs can borrow depends on their observable characteristics, and the entrepreneurs' assets act as collateral for their debts. Since the implicit rate of return for entrepreneurs is higher than the rate for workers, entrepreneurs have a higher saving rate, which is consistent with the data. We calibrate the parameters of the model to match key moments of the data and discuss the implications of the model and its components for entrepreneurial choice and wealth inequality. We show that our model with entrepreneurial choice matches very well the observed distribution of wealth, for both entrepreneurs and nonentrepreneurs.

Section~\ref{sec:microlit} first documents the relationship between wealth and entrepreneurship and then surveys the evidence that entrepreneurs are borrowing constrained. Section~\ref{sec:model} describes the model and our calibration procedure. Section~\ref{sec:comparing} discusses the role of entrepreneurship and voluntary bequests in generating large wealth concentration and studies the aggregate effects of changing the borrowing constraints. Section~\ref{sec:hank} inspects further the mechanisms at work in our model and compares their observable implications to those in the observed data. Section~\ref{sec:conclusion} presents conclusions.

\smartbib

\end{document}
